\documentclass[11pt,a4paper]{article}
% Common preamble for the hint sheets (\input from each file)
\usepackage[utf8]{inputenc}
\usepackage[T1]{fontenc}
\usepackage{lmodern}
\usepackage{amsmath,amssymb,amsthm,mathtools}
\usepackage[a4paper,margin=2.3cm]{geometry}
\usepackage{enumitem}
\usepackage{xcolor}
\usepackage[hidelinks]{hyperref}
\renewcommand{\leq}{\leqslant}
\renewcommand{\le}{\leqslant}
\renewcommand{\geq}{\geqslant}
\renewcommand{\ge}{\geqslant}
\setlength{\parindent}{0pt}
\setlength{\parskip}{0.35em}
% \hintstitle{Sheet title}{Class work / Homework}
\newcommand{\hintstitle}[2]{\begin{center}{\LARGE\bfseries Hints: #1\par}\vspace{0.5em}{\large #2\par}\end{center}
\noindent\textit{The hints to each problem are cumulative. Read one, try again, and open the next only if you are still stuck.}\par\medskip\hrule\medskip}
\newcommand{\problem}[1]{\par\bigskip{\large\bfseries Problem #1}\par\nopagebreak\smallskip}
\newcommand{\subpart}[1]{\par\smallskip\textbf{(#1)}\par\nopagebreak}
\newcommand{\hint}[1]{\par\smallskip\noindent\textcolor{black!70}{\textbf{Hint #1.}}\enspace}
\newcommand{\note}[1]{\par\smallskip\noindent\textit{Note.} #1\par}
\providecommand{\lt}{<}
\providecommand{\gt}{>}
\providecommand{\tightlist}{\setlength{\itemsep}{0pt}\setlength{\parskip}{0pt}}
\emergencystretch=3em

\begin{document}
\hintstitle{Number theory -- 6}{Class work}

\problem{1}
\hint{1} Write $(f*g)(n)=\sum_{ab=n}f(a)g(b)$. Then commutativity is obvious, and associativity holds because both sides equal $\sum_{abc=n}f(a)g(b)h(c)$.
\hint{2} Inverse: solve $(f*g)(n)=\varepsilon(n)$ recursively. Set $g(1)=1/f(1)$, and for $n>1$ put $g(n)=-\frac1{f(1)}\sum_{d\mid n,\,d<n}f(n/d)g(d)$. Multiplicativity: for coprime $m,n$, the divisors of $mn$ are exactly the products $d_1d_2$ with $d_1\mid m$, $d_2\mid n$.
\hint{3} Inverses: $1^{-1}=\mu$ (Theorem 2). Since $d=1*1$, $d^{-1}=\mu*\mu$, which is multiplicative with $p\mapsto-2$, $p^2\mapsto1$, $p^k\mapsto0$ for $k\ge3$. Since $\sigma=\mathrm{id}*1$ and $\mathrm{id}^{-1}=\mu\cdot\mathrm{id}$, $\sigma^{-1}=(\mu\,\mathrm{id})*\mu$, with $p\mapsto-(p+1)$, $p^2\mapsto p$, $p^k\mapsto0$ for $k\ge3$.

\problem{2}
\hint{1} Grouping: the fractions $\frac kn$, $1\le k\le n$, written in lowest terms, have denominators $d\mid n$. Each denominator $d$ occurs with exactly $\varphi(d)$ numerators. Multiplicativity: both sides of $\sum_{d\mid n}\varphi(d)=n$ are multiplicative, and at $p^k$ the left side telescopes.
\hint{2} For $\sum_{d\mid n}\mu(d)=[n=1]$, use $(1-1)^{\omega(n)}$, or multiplicativity. Then $g=f*1\iff f=g*\mu$ by associativity. From $\varphi*1=\mathrm{id}$ you get $\varphi=\mu*\mathrm{id}$, i.e.\ $\varphi(n)=\sum_{d\mid n}\mu(d)\frac nd=n\prod_{p\mid n}(1-\frac1p)$.
\hint{3} The last two identities are $\mu*d=\mu*1*1=1$ and $\mu*\sigma=\mu*1*\mathrm{id}=\mathrm{id}$.

\problem{3}
\hint{1} Compute small cases. Primes work: $p(p+1)\equiv1\cdot2\pmod{p-1}$. Also $n=1,4,6,22$ work. The task is to show nothing else does.
\hint{2} If $p^2\mid n$, then $p\mid\varphi(n)$ and $p\mid n\sigma(n)$, so $p\mid2$. Hence $n=2^a m$ with $m$ odd and squarefree. If $a\ge3$, then $4\mid\varphi(n)$ and $4\mid n$, which is impossible. So $a\le2$.
\hint{3} If $m$ has two odd prime factors, then $4\mid\varphi(n)$. But $\sigma(m)=\prod(p+1)$ is divisible by $4$, so $n\sigma(n)\equiv0\not\equiv2\pmod4$.
\hint{4} The remaining cases are $n=2,4$, $n=2p$ and $n=4p$. For $n=2p$, reducing mod $p-1$ gives $p-1\mid10$. For $n=4p$, it gives $p-1\mid54$. Check the few candidates. Answer: $n=1,4,6,22$ and all primes.

\problem{4}
\hint{1} $\sum_{n\le x}d(n)$ counts lattice points $(a,b)$ with $ab\le x$. By the symmetry $a\leftrightarrow b$ it equals $2\sum_{a\le\sqrt x}\lfloor x/a\rfloor-\lfloor\sqrt x\rfloor^2$.
\hint{2} Insert $\lfloor x/a\rfloor=x/a+O(1)$ and $\sum_{a\le y}\frac1a=\ln y+\gamma+O(1/y)$. All error terms are $O(\sqrt x)$, and the main terms give $x\ln x+(2\gamma-1)x$.
\hint{3} Coprime pairs: by Möbius inversion, the number of coprime pairs in $\{1,\dots,N\}^2$ is $\sum_{d\le N}\mu(d)\lfloor N/d\rfloor^2=N^2\sum_{d\le N}\frac{\mu(d)}{d^2}+O(N\ln N)$.
\hint{4} The tail of $\sum\mu(d)/d^2$ is $O(1/N)$, so the probability tends to $\sum_{d\ge1}\mu(d)d^{-2}=1/\zeta(2)$. Only here does the value $\zeta(2)=\pi^2/6$ enter (Euler's evaluation, from Series); the rest uses only $\mu*1=\varepsilon$.

\problem{5}
\hint{1} Euclid: if $2^k-1$ is prime, then $\sigma(2^{k-1}(2^k-1))=(2^k-1)2^k$ by multiplicativity.
\hint{2} Euler: write $n=2^{k-1}m$ with $m$ odd and $k\ge2$. Then $(2^k-1)\sigma(m)=2^km$, so $m=(2^k-1)t$ and $\sigma(m)=2^kt=m+t$. Since $t<m$ are both divisors of $m$, the divisors of $m$ are only $t$ and $m$. Hence $t=1$ and $m$ is prime.
\hint{3} Odd perfect $n$: for a prime power, $\frac{\sigma(p^k)}{p^k}<\frac p{p-1}\le\frac32<2$. With two distinct odd primes, $\frac{\sigma(n)}n<\frac p{p-1}\cdot\frac q{q-1}\le\frac32\cdot\frac54=\frac{15}8<2$.
\end{document}
